\section{Results}

\subsection{Main Results}

Table~\ref{tab:firm-size-wage-premium-regressions} reports the regression results from the main specification in Equation (\ref{eq:main}). Column (1) shows the raw log-wage gradient by firm size. The coefficient for firms with 101 or more workers is 1.019 log points, consistent with the large raw wage premium shown in the descriptive statistics. Column (2) adds sector-year fixed effects, reducing the coefficient for 101+ firms to 0.819 log points. Column (3) adds gender, age, age squared, and education controls, lowering the coefficient to 0.598. Column (4) adds formality, reducing the coefficient further to 0.359. Thus, much of the raw firm-size wage gap is associated with sector-year conditions and observable differences in worker and job characteristics, including gender, age, education, and formality. Nonetheless, a positive and statistically significant conditional premium remains after controlling for these observables, suggesting that barriers to firm growth may matter for workers' wages. If larger firms are more productive or better able to share rents with workers, then obstacles that prevent small firms from scaling up, or workers from moving into larger firms, can depress earnings and contribute to wage inequality and labor-market inefficiencies.

\begin{table}[htbp]
  \centering
  \caption{Firm-size wage premium regressions: the role of formality}
  \label{tab:firm-size-wage-premium-regressions}
  \small
  \begin{tabular}{lcccc}
    \toprule
    & \multicolumn{4}{c}{Dependent variable: log real hourly labor income} \\
    \cmidrule(lr){2-5}
    & (1) & (2) & (3) & (4) \\
    \midrule
    Firm size: 2-3 & 0.139*** & 0.193*** & 0.168*** & 0.145*** \\
     & (0.033) & (0.029) & (0.027) & (0.029) \\
    Firm size: 4-5 & 0.287*** & 0.325*** & 0.293*** & 0.234*** \\
     & (0.036) & (0.034) & (0.035) & (0.044) \\
    Firm size: 6-10 & 0.390*** & 0.402*** & 0.337*** & 0.234*** \\
     & (0.034) & (0.038) & (0.041) & (0.052) \\
    Firm size: 11-19 & 0.492*** & 0.468*** & 0.363*** & 0.224*** \\
     & (0.039) & (0.043) & (0.044) & (0.050) \\
    Firm size: 20-30 & 0.529*** & 0.479*** & 0.356*** & 0.190*** \\
     & (0.042) & (0.047) & (0.046) & (0.050) \\
    Firm size: 31-50 & 0.592*** & 0.521*** & 0.378*** & 0.195*** \\
     & (0.052) & (0.051) & (0.048) & (0.049) \\
    Firm size: 51-100 & 0.657*** & 0.567*** & 0.407*** & 0.215*** \\
     & (0.061) & (0.053) & (0.047) & (0.047) \\
    Firm size: 101+ & 0.843*** & 0.672*** & 0.480*** & 0.282*** \\
     & (0.094) & (0.069) & (0.063) & (0.064) \\
    \midrule
    Gender, age, and education controls & No & No & Yes & Yes \\
    Formality control & No & No & No & Yes \\
    Sector-year fixed effects & No & Yes & Yes & Yes \\
    Observations & 3,852,322 & 3,852,322 & 3,852,322 & 3,852,322 \\
    $R^2$ & 0.162 & 0.221 & 0.353 & 0.362 \\
    \bottomrule
  \end{tabular}
  \vspace{0.3em}
  \begin{minipage}{0.95\textwidth}
  \footnotesize
  Notes: The omitted category is solo workers. All columns use the same estimation sample and GEIH expansion weights. Standard errors, clustered by sector, are reported in parentheses. Worker controls include a female-worker indicator, age, age squared, and education dummies. Column (4) adds labor formality. Significance levels: * $p<0.10$, ** $p<0.05$, *** $p<0.01$.
  \end{minipage}
\end{table}


%The comparison across columns is central to the interpretation. The raw premium documents the inequality that workers actually face across the firm-size distribution. The reduction after adding education, age, gender, sector-year fixed effects, and formality shows that part of the gap reflects sorting and segmentation: larger firms employ a more formal and more educated workforce and operate in different sector-year environments. The remaining conditional premium does not prove that larger firms causally raise wages, but it is consistent with firm-level wage-setting differences, rent sharing, internal labor markets, technology, productivity, or unobserved worker-firm matching. In a labor market where most workers are concentrated in solo work and micro-firms, this remaining premium is economically meaningful even after the controls absorb a large share of the raw gap.

%Figure~\ref{fig:firm-size-premium-attenuation} summarizes this attenuation for selected firm-size categories. The figure shows that the raw gap is large, but that sector-year conditions, worker composition, and especially formality absorb an important share of the difference between small and large firms.

%\begin{figure}[htbp]
%  \centering
%  \includegraphics[width=0.9\textwidth]{fig68_\figurelang.png}
%  \caption{Attenuation of the firm-size wage premium across specifications}
%  \label{fig:firm-size-premium-attenuation}
%\end{figure}

Figure~\ref{fig:firm-size-wage-premium} transforms the coefficients from Column (4) in Table (\ref{tab:firm-size-wage-premium-regressions}) into percentage premiums. Relative to solo workers, the conditional premium is 20.7 percent for firms with 2--3 workers, 36.0 percent for firms with 4--5 workers, 37.3 percent for firms with 6--10 workers, roughly 31--35
percent for intermediate size categories, and 43.2 percent for firms with 101 or more workers. All confidence intervals exclude zero.

\begin{figure}[htbp]
  \centering
  \includegraphics[width=0.9\textwidth]{fig57_\figurelang.png}
  \caption{Conditional firm-size wage premium relative to solo workers}
  \label{fig:firm-size-wage-premium}
\end{figure}

\subsection{Dynamic results}

Figure~\ref{fig:firm-size-wage-premium-over-time} estimates the full specification separately by interacting firm-size categories with year indicators. The conditional premium is positive throughout the sample period in all firm-size categories. The largest firms display the highest premiums in most years, although the confidence intervals are wider for these categories.

\begin{figure}[htbp]
  \centering
  \includegraphics[width=0.95\textwidth]{fig59_\figurelang.png}
  \caption{Conditional firm-size wage premium over time}
  \label{fig:firm-size-wage-premium-over-time}
\end{figure}

Table~\ref{tab:firm-size-premium-trend-test} formally tests whether the premium changed over time using a linear trend interacted with firm-size categories. The estimated trend coefficients are positive for all size categories. The implied increase between 2008 and 2025 ranges from 3.2 percent for firms with 20--30 workers to 7.9 percent for firms with 101 or more workers. However, the two-sided tests do not reject a zero trend at the 5 percent level. Thus, the current evidence supports a positive conditional firm-size wage premium, but the evidence that this premium has changed over time is suggestive rather than statistically conclusive. Since all estimated trend coefficients are positive, the suggestive pattern points toward widening rather than narrowing, but the bilateral tests do not support a statistically conclusive trend.

The dynamic exercise also clarifies what it would mean for the firm-size wage gradient to steepen or flatten over time. A steepening premium would imply that access to large-firm employment has become increasingly important for wages, raising concerns about inequality and the allocation of workers across the productive structure. A flattening premium would be ambiguous: it could reflect improving wages and productivity among smaller firms, but it could also reflect weaker wage growth among larger firms. For this reason, the policy-relevant
object is not only the gap between small and large firms, but the joint evolution of wage levels, employment shares, formality, and productivity across firm-size categories.

\begin{table}[htbp]
  \centering
  \caption{Testing whether the firm-size wage premium changed over time}
  \label{tab:firm-size-premium-trend-test}
  \small
  \begin{tabular}{lcccc}
    \toprule
    Firm size & Annual trend & S.E. & $p$-value & Implied 2008--2025 change \\
    \midrule
    2-3 & 0.0024 & (0.0029) & 0.409 & 4.3\% \\
    4-5 & 0.0005 & (0.0039) & 0.909 & 0.8\% \\
    6-10 & 0.0014 & (0.0036) & 0.693 & 2.5\% \\
    11-19 & 0.0022 & (0.0043) & 0.610 & 3.9\% \\
    20-30 & 0.0023 & (0.0042) & 0.595 & 3.9\% \\
    31-50 & 0.0029 & (0.0046) & 0.533 & 5.1\% \\
    51-100 & 0.0037 & (0.0042) & 0.393 & 6.5\% \\
    101+ & 0.0054 & (0.0047) & 0.268 & 9.7\% \\
    \bottomrule
  \end{tabular}
  \vspace{0.3em}
  \begin{minipage}{0.95\textwidth}
  \footnotesize
  Notes: The annual trend is the coefficient on the interaction between firm-size category and a linear year trend, with solo workers as the omitted category. The specification controls for gender, age, age squared, education, formality, and sector-year fixed effects, and uses GEIH expansion weights. Standard errors are clustered by sector. The $p$-value corresponds to the two-sided test that the trend differs from zero. The final column reports $100\times[\exp(17\hat{\delta})-1]$, the implied change in the firm-size wage ratio between 2008 and 2025.
  \end{minipage}
\end{table}


\subsection{Formality non-linear results}

Figure~\ref{fig:firm-size-formality-premium} asks whether the nonlinear pattern seen in Figure~\ref{fig:income-formality-comparison} survives after controlling for worker composition and sector-year conditions. The figure reports firm-size premiums relative to solo workers within each formality group. Among informal workers, the conditional premium is positive and statistically significant in every firm-size category: it is 21.7 percent in firms with 2--3 workers, rises to about 48--51 percent across most categories from 6 to 100 workers, and remains 40.0 percent in firms with 101 or more workers. Among formal workers, the within-formality gradient is much flatter and less precisely estimated: the estimated premium ranges from 2.1 to 7.0 percent for most categories and is 16.5 percent for firms with 101 or more workers, but the confidence intervals include zero throughout when each category is compared with formal solo workers. This does not mean, however, that there is no wage structure within formal employment. The U-shape in Figure~\ref{fig:income-formality-comparison} points to a different comparison: formal workers in 101+ firms versus formal workers in intermediate-size firms. Direct contrasts from the interacted model show that formal workers in 101+ firms earn between 8.9 and 14.2 percent more than formal workers in every smaller non-solo category, and all of these contrasts are statistically significant. Table~\ref{tab:formal-101-contrasts} reports these contrasts explicitly. Thus, the estimates are consistent with a heterogeneous and selected formal solo group coexisting with a large-firm premium relative to
the middle of the formal firm-size distribution. The informal sector, by contrast, displays a robust conditional firm-size gradient relative to solo work.

\begin{figure}[htbp]
  \centering
  \includegraphics[width=0.95\textwidth]{fig74_\figurelang.png}
  \caption{Adjusted firm-size wage premium by formality status}
  \label{fig:firm-size-formality-premium}
\end{figure}

Figure~\ref{fig:adjusted-formality-endpoints} provides the closest
regression-adjusted counterpart to
Figure~\ref{fig:income-formality-comparison}. The y-axis is the adjusted percent
premium relative to solo workers within the same formality-year group, rather
than an adjusted wage level. This normalization avoids choosing arbitrary
covariate values and makes the comparison with the solo-worker benchmark
explicit. Among formal workers, the raw U-shape largely disappears after
controls: in 2025, the adjusted estimates for all categories below 101+ range
from -3.7 to 0.6 percent, and the 101+ estimate is 10.5 percent with a
confidence interval that includes zero. The 2008 formal profile is similar.
Among informal workers, by contrast, the adjusted gradient remains steep and
statistically significant: in 2025, the premium rises from 24.3 percent in firms
with 2--3 workers to 53.6 percent in firms with 101 or more workers. Thus, the
descriptive U-shape among formal workers appears to be driven mainly by worker
composition within formal solo work, while the
informal firm-size premium survives the controls. This conclusion concerns the
comparison with solo workers; Table~\ref{tab:formal-101-contrasts} turns to the
separate comparison between large formal firms and the middle of the formal
firm-size distribution.

\begin{figure}[htbp]
  \centering
  \includegraphics[width=0.95\textwidth]{fig76_\figurelang.png}
  \caption{Adjusted firm-size wage premium by formality status, 2008 and 2025}
  \label{fig:adjusted-formality-endpoints}
\end{figure}

\begin{table}[htbp]
  \centering
  \caption{Formal-worker large-firm wage premium relative to other formal firm-size categories}
  \label{tab:formal-101-contrasts}
  \small
  \begin{tabular}{lccc}
    \toprule
    Reference category & Premium & 95\% CI & $p$-value \\
    \midrule
    Solo workers & 12.4\% & [-6.5, 35.1] & 0.213 \\
    2-3 & 12.5\% & [4.5, 21.0] & 0.002 \\
    4-5 & 12.9\% & [6.3, 19.9] & $<0.001$ \\
    6-10 & 13.1\% & [6.1, 20.5] & $<0.001$ \\
    11-19 & 11.3\% & [2.7, 20.6] & 0.009 \\
    20-30 & 13.3\% & [4.4, 22.9] & 0.003 \\
    31-50 & 11.2\% & [3.3, 19.7] & 0.005 \\
    51-100 & 8.1\% & [1.6, 15.0] & 0.013 \\
    \bottomrule
  \end{tabular}
  \vspace{0.3em}
  \begin{minipage}{0.92\textwidth}
  \footnotesize
  Notes: Each row compares formal workers in firms with 101 or more workers with formal workers in the reference firm-size category. Estimates come from the firm-size-by-formality specification with gender, age, age squared, education, and sector-year fixed effects. Premiums are computed as $100\times[\exp(\hat{\beta})-1]$ from the relevant linear contrast.
  \end{minipage}
\end{table}


Figure~\ref{fig:large-firm-formality-premium-over-time} examines whether the
pooled formal-worker result masks time variation. The figure plots the adjusted
premium for firms with 101 or more workers relative to solo workers, separately
within formal and informal employment in each year. The evidence does not show a
reappearance of a statistically significant large-firm premium among formal
workers in the most recent years. The formal 101+ premium is positive and
significant in several earlier years, especially 2009--2013 and 2017, but after
2021 the point estimates are smaller and the confidence intervals include zero;
in 2025, the estimate is 10.5 percent with a confidence interval from -8.2 to
33.0 percent. By contrast, the informal 101+ premium is positive and significant
throughout the period, rising from 39.4 percent in 2008 to 53.6 percent in
2025.

\begin{figure}[htbp]
  \centering
  \includegraphics[width=0.95\textwidth]{fig75_\figurelang.png}
  \caption{Adjusted large-firm wage premium by formality status over time; 2020 is omitted}
  \label{fig:large-firm-formality-premium-over-time}
\end{figure}



\subsection{Interpretation: Inequality, Efficiency, and Misallocation}

The results connect the firm-size wage premium to three related issues. The
first is wage inequality. Since workers in larger firms earn more, the allocation
of workers across firm sizes affects the overall wage distribution. The second is
labor-market segmentation. The strong role of formality in reducing the
coefficient for large firms suggests that the firm-size gradient is partly a
formal-informal gradient. The third is efficiency. If larger firms are more
productive and the wage premium partly reflects higher marginal products or
rents shared with workers, then the large employment mass in solo work and
micro-firms may indicate that workers and firms are not moving freely toward
higher-productivity matches.

This last interpretation should be treated as a hypothesis rather than as a
direct conclusion. The GEIH does not observe firms, capital, value added, or
marginal products, so the estimates cannot establish misallocation in the strict
sense of \citet{hsieh2009misallocation}. What the paper documents is a set of
facts consistent with a misallocation concern: employment is heavily concentrated
in very small units, larger firms pay substantially more, and a positive premium
remains after conditioning on observed worker characteristics, formality, and
sector-year fixed effects. Establishing the productivity channel directly would
require matched employer-employee data or firm-level information on output,
capital, and employment.

\subsection{Limitations}

Several limitations qualify the interpretation. First, firm size is reported by
workers in a household survey, and the data do not identify employers. Second,
the analysis cannot separate worker sorting from firm wage-setting: a remaining
conditional premium may reflect unobserved worker ability, firm rents, internal
labor markets, productivity differences, or match quality. Third, formality is
both a control and a potential mechanism. Adding it to the regression is useful
because it shows how much of the raw gradient is associated with the formal
segment of the labor market, but it may also absorb part of the channel through
which larger firms pay more. Finally, because the GEIH does not observe value
added, capital, or marginal products, the evidence should be read as suggestive
of allocative frictions rather than as a direct test of misallocation.
